\subsection{Bestimmung des Arbeitspunkts und der Leistung einer Diode}
% Alt: Aufgabe 4

\Aufgabe{
    \label{Aufg4}
    Eine Diode hat die folgende Kennlinie im
    Durchlassbereich. Sie wird in Reihe mit einer
    idealen Spannungsquelle $U_0=1\,\mathrm{V}$ und einem
    Widerstand mit dem Wert von $R=200\,\Omega$ geschaltet.

    \begin{figure}[H]
        \centering
        \begin{subfigure}[h]{0.45\textwidth}
            \centering
            \input{Tikz/src/FigDiodeAP}\alttext{Zu sehen sind zwei Abbildungen. Links ist eine Reihenschaltung mit drei Bauteilen dargestellt:
                Eine Spannungsquelle, bezeichnet mit \(U_0\), deren Pfeil nach unten zeigt, ein Widerstand, gekennzeichnet mit dem Buchstaben \(R\) und eine Diode, symbolisiert durch das typische Dreieck mit einem Strich an der Spitze und beschriftet mit \(D\).
                Der Stromfluss in der Schaltung wird durch einen roten Pfeil markiert, der im Uhrzeigersinn vom Widerstand zur Diode verläuft und mit \(I\) bezeichnet ist. Neben der Diode befindet sich ein weiterer Spannungspfeil in Richtung des Stromflusses, beschriftet mit \(U_\mathrm{D}\).}
        \end{subfigure}
        \begin{subfigure}[h]{0.45\textwidth}
            \centering
            \includesvg[width=0.9\textwidth]{Bilder/Aufgaben/FigDiodeAP}\alttext{Rechts daneben befindet sich ein Diagramm, das den Zusammenhang zwischen Stromstärke \(I\) und Spannung U über der Diode darstellt. Die horizontale Achse zeigt die Spannung in Volt (V), von 0 bis etwas über 1 Volt. Die vertikale Achse zeigt den Strom in Milliampere (mA), von \(0\) bis \(7 \ mA\).
                Die Kurve im Diagramm beginnt nahe null bei niedrigen Spannungen und steigt ab etwa 0,6 Volt steil an - dies verdeutlicht das typische Verhalten einer Siliziumdiode: bei einer bestimmten Schwellenspannung fließt plötzlich deutlich mehr Strom.}
        \end{subfigure}
        \label{fig:FigDiodeAP}
    \end{figure}


    \begin{itemize}
        \item[a)]
              Welcher Strom $I$ und welche Spannung $U_\mathrm{D}$
              stellen sich ein? Lösen Sie das Problem graphisch.
        \item[b)]
              Welche elektrische Leistung wird an der Diode umgesetzt?
        \item[c)]
              Wie muss $R$ gewählt werden, damit sich ein Strom von $3\,\mathrm{mA}$ einstellt?
    \end{itemize}
}

\Loesung{
    \begin{itemize}
        \item[a)]
              Im ersten Schritt wird die Wirkung der Spannungsquelle auf den Widerstand $R$ betrachtet.
              Bei $U_0=1\,\mathrm{V}$ fließen durch den Widerstand $R$ ein Strom von $I_\mathrm{R}=\frac{U_0}{R}=\frac{1\,\mathrm{V}}{200\,\Omega}=5\,\mathrm{mA}$.
              Es kann die Widerstandgerade vom Ursprung aus zu dem Punkt ($1\,\mathrm{V}$, $5\,\mathrm{mA}$) eingezeichnet werden (linke Abbildung).\\
              Gemäß des Aufbaus der Schaltung muss die Widerstandsgerade im nächsten Schritt gespiegelt werden, es entsteht die rechte Darstellung.
              Der Schnittpunkt beider Linien entspricht dem Arbeitspunkt ($U_\text{AP}$, $I_\text{AP}$), dieser lautet: ($0,75\,\mathrm{V}$, $1,2\,\mathrm{mA}$).

              \begin{minipage}{\linewidth}
                  \begin{figure}[H]
                      \centering
                      \begin{subfigure}[b]{0.49\linewidth}
                          \centering
                          \includesvg[width=0.9\textwidth]{Bilder/Aufgaben/LsgDiodeAP1}
                          \alttext{Die Abbildung zeigt zwei Diagramme zur Bestimmung des Arbeitspunktes einer Diode.
                          Auf der x-Achse ist die Diodenspannung \(U_\mathrm{D}\) in Volt und auf der y-Achse der Strom \(I\) in Milliampere aufgetragen. Die blaue Kurve stellt die nichtlineare Kennlinie der Diode \(I_\mathrm{D}\)dar. Der schwarze Graph zeigt jeweils die Lastgerade \(I_\mathrm{R}\),die sich aus dem angeschlossenen Widerstand ergibt.
                          Der Schnittpunkt liegt nahe dem steilen Bereich der Kennlinie.}
                      \end{subfigure}
                      \begin{subfigure}[b]{0.49\linewidth}
                          \centering
                          \includesvg[width=0.9\textwidth]{Bilder/Aufgaben/LsgDiodeAP2}
                          \alttext{Das rechte Diagramm stellt ebenfalls auf der x-Achse die Diodenspannung \(U_\mathrm{D}\) in Volt und auf der y-Achse den Strom \(I\) in Milliampere dar. Die blaue Kurve stellt die nichtlineare Kennlinie der Diode \(I_\mathrm{D}\) dar. Der schwarze Graph zeigt jeweils die Lastgerade \(I_\mathrm{R}\).
                          Die Lastgerade und die Kennlinie der Diode schneiden sich im Punkt \(0,75 \ \mathrm{V}\) und \(1,2 \ \mathrm{mA}\). Dieser Schnittpunkt ist markiert und kennzeichnet den tatsächlichen Arbeitspunkt.}
                      \end{subfigure}
                  \end{figure}
              \end{minipage}

        \item[b)]
              Die Leistung ergibt sich aus dem zuvor ermittelten Arbeitspunkt.
              \begin{equation*}
                  P = U_\text{AP} \cdot I_\text{AP} = 0,75 \, \text{V} \cdot 1,2 \, \text{mA} = 0,9 \, \text{mW}
              \end{equation*}

        \item[c)]
              Ausgehend vom Punkt ($1\,\mathrm{V}, 0\,\mathrm{mA}$) zum Schnittpunkt mit der Diodenkennlinie bei $I=3\,\mathrm{mA}$ ergibt sich die gesuchte Widerstandsgerade.

              \begin{minipage}{\linewidth}
                  \begin{figure}[H]
                      \centering
                      \includesvg[width=0.45\linewidth]{Bilder/Aufgaben/LsgDiodeAP3}
                      \alttext{Die Abbildung zeigt das Diagramm der Diodenkennlinie. Auf der x-Achse ist die Diodenspannung \(U_\mathrm{D}\) in Volt und auf der y-Achse der Strom \(I\) in Milliampere aufgetragen. Die blaue Kurve stellt die nichtlineare Kennlinie der Diode \(I_\mathrm{D}\)dar.
                      Die Lastgerade \(I_\mathrm{R}\) startet bei \(1\,\mathrm{V}, 0\,\mathrm{mA}\) und schneidet die Diodenkennlinie \(I_\mathrm{D}\) bei \(0,86 \,\mathrm{V}, 3\,\mathrm{mA}\). Der Schnittpunkt ist mit einem Kreis umrundet.}
                  \end{figure}
              \end{minipage}

              Es kann der Schnittpunkt ($0,86\,\mathrm{V}, 3\,\mathrm{mA}$) abegelesen werden.

              \begin{gather*}
                  U_\mathrm{R} = U_0 - U_\mathrm{D} = 1 \, \mathrm{V} - 0,86 \, \mathrm{V} = 0,14 \, \mathrm{V} \\
                  I = 3 \, \mathrm{mA} \\
                  R = \frac{U_{\text{R}}}{I} = \frac{0,14 \, \text{V}}{3 \, \text{mA}} = 46,67 \, \Omega
              \end{gather*}
    \end{itemize}
    Siehe: Abschnitt \ref{fig:Diodenkennlinie}
}